\section*{第 \ref{ch:g10:synthesis-yield} 章——合成：产率与纯度}

\begin{solution}{exo:g10:synthesis-yield:1}
$\eta = 4.2 / 5.0 = 0.84$，产率为 \qty{84}{\percent}。
\end{solution}

\begin{solution}{exo:g10:synthesis-yield:2}
称量反应物；加热回流；对粗产物进行减压过滤；
重结晶；测量熔点。
\end{solution}

\begin{solution}{exo:g10:synthesis-yield:3}
水从底部进入，先充满整个夹套，
再从顶部流出，因此冷凝器的全长都得到冷却。顶端
保持敞开，使装置不密闭：加热密闭装置
会使压强升高，可能导致装置破裂。
\end{solution}

\begin{solution}{exo:g10:synthesis-yield:4}
产物须在热溶剂中极易溶解，
而在冷溶剂中几乎不溶。杂质含量很少，
会留在冷溶剂中，最终进入滤液。
\end{solution}

\begin{solution}{exo:g10:synthesis-yield:5}
不能。它在低于 \qty{135}{\celsius} 时即开始熔化，且熔化范围跨越数度，
这是固体不纯的表现。应重结晶、干燥，
然后重新测定熔点。
\end{solution}

\begin{solution}{exo:g10:synthesis-yield:6}
$n = 2.00 / 138.0 = \qty{0.0145}{mol}$ 水杨酸，为限量
反应物；$m_{\max} = 0.0145 \times 180.0 = \qty{2.61}{g}$；
$\eta = 2.03 / 2.61 = 0.78$，即 \qty{78}{\percent}。
\end{solution}

\begin{solution}{exo:g10:synthesis-yield:7}
称得的质量包含水分，因此比阿司匹林的质量大，
计算产率偏高，这里甚至高于
\qty{100}{\percent}，显然不合理。干燥的粗产物仍含杂质
（未反应的水杨酸、副产物），它们同样增加质量，因此其产率
与纯阿司匹林的真实产率相比同样偏高。
\end{solution}

\begin{solution}{exo:g10:synthesis-yield:8}
减压过滤快得多，且晶体干燥得多，因为抽气能
抽出大部分液体；晶体也便于在漏斗上洗涤，
并便于从平整的滤纸上刮下。
\end{solution}

\begin{solution}{exo:g10:synthesis-yield:9}
产物 $3.1/5.0 = 0.62$；阿司匹林 $3.1/5.0 = 0.62$；水杨酸
$2.2/5.0 = 0.44$。产物的迁移行为与阿司匹林相同，且未出现
水杨酸的斑点：可知产物为阿司匹林，未检出水杨酸。
\end{solution}

\begin{solution}{exo:g10:synthesis-yield:10}
溶剂 A：热时极易溶解，冷时几乎不溶；溶剂 B 即使在冷时，也使大部分
产物保持溶解。使用 A 时，\qty{3.0}{g} 约需
$3.0 / 150 = \qty{0.020}{L}$，即 \qty{20}{mL} 热溶剂；
冷却后，最多有 $2 \times 0.020 = \qty{0.04}{g}$ 保持溶解。
\end{solution}

\begin{solution}{exo:g10:synthesis-yield:11}
沸腾的液体上方需留有空间：泡沫和飞溅物不得到达
冷凝器。沸石为蒸气提供形成气泡的微小位置，
使液体平稳沸腾，而非突然爆沸。
\end{solution}

\begin{solution}{exo:g10:synthesis-yield:12}
$0.80 \times 0.70 = 0.56$：\qty{56}{\percent}。三步、每步
\qty{90}{\percent}：$0.90^3 = 0.73$，\qty{73}{\percent}。每一步
都会损失产物（并损失时间与金钱），且损失逐步累积：步骤越少，
最终剩余的产物越多。
\end{solution}

\begin{solution}{exo:g10:synthesis-yield:13}
最多可得 $3.3 \times 0.050 = \qty{0.17}{g}$ 阿司匹林。冰冷的水
溶解的阿司匹林更少，因为其溶解度随
温度降低而下降。
\end{solution}

\begin{solution}{exo:g10:synthesis-yield:14}
\begin{enumerate}
  \item $n_0(\text{水杨酸}) = 2.76 / 138.0 = \qty{0.0200}{mol}$，
        乙酸酐 \qty{0.0500}{mol}：水杨酸为限量反应物，
        $x_{\max} = \qty{0.0200}{mol}$，$0.0500 - 0.0200 =
        \qty{0.0300}{mol}$ 乙酸酐剩余。
  \item 合成生成 \qty{0.0200}{mol} 乙酸，
        过量部分 $2 \times 0.0300 = \qty{0.0600}{mol}$ 被分解，
        共计 \qty{0.0800}{mol}。
\end{enumerate}
\end{solution}

\begin{solution}{exo:g10:synthesis-yield:15}
$n(\text{阿司匹林}) = \num{1.00e6} / 180.0 = \qty{5.56e3}{mol}$。产率为
\qty{85}{\percent} 时，所需水杨酸为：
$\num{5.56e3} / 0.85 = \qty{6.54e3}{mol}$，即
$\num{6.54e3} \times 138.0 = \qty{9.0e5}{g}$，约 \qty{0.90}{t}。
\end{solution}

\begin{solution}{pb:g10:synthesis-yield:1}
\textbf{1.} C：$7 + 4 = 11$ 与 $9 + 2 = 11$；H：$6 + 6 = 12$ 与
$8 + 4 = 12$；O：$3 + 3 = 6$ 与 $4 + 2 = 6$。已配平。

\textbf{2.} 水杨酸 $7 \times 12.0 + 6 \times 1.0 + 3 \times 16.0
= \qty{138.0}{g/mol}$；乙酸酐 $4 \times 12.0 + 6 \times 1.0 +
3 \times 16.0 = \qty{102.0}{g/mol}$；阿司匹林 $9 \times 12.0 + 8 \times
1.0 + 4 \times 16.0 = \qty{180.0}{g/mol}$；乙酸 $2 \times 12.0
+ 4 \times 1.0 + 2 \times 16.0 = \qty{60.0}{g/mol}$。

\textbf{3.} 反应物为水杨酸和乙酸酐；所需产物为
阿司匹林；副产物为乙酸。

\textbf{4.} 加热可加快反应；加热回流时，蒸气
经冷凝后回流，因此反应物和产物均不会损失。

\textbf{5.} $n = 3.0 / 138.0 = \qty{0.0217}{mol}$。

\textbf{6.} $m = 6.0 \times 1.08 = \qty{6.48}{g}$；
$n = 6.48 / 102.0 = \qty{0.0635}{mol}$。

\textbf{7.} 水杨酸：$0.0217 - x$；乙酸酐 $0.0635 - x$；
阿司匹林和乙酸：$x$。水杨酸先耗尽：
$x_{\max} = \qty{0.0217}{mol}$，水杨酸为限量反应物。

\textbf{8.} $m_{\max} = 0.0217 \times 180.0 = \qty{3.91}{g}$。

\textbf{9.} 为使水杨酸全部转化为阿司匹林，若有剩余，
会与阿司匹林混在一起而难以除去，
而过量的乙酸酐会在最后加入的水中被分解，
其生成的乙酸则随洗涤除去。

\textbf{10.} 产物含水：若以 \qty{3.6}{g} 计，则会得出虚假的
$3.6 / 3.91 = \qty{92}{\percent}$。

\textbf{11.} $3.3 / 3.91 = 0.84$：\qty{84}{\percent} 为粗
干产物，且不纯。

\textbf{12.} 杂质被除去，但部分阿司匹林
仍溶于冷的溶剂和洗涤液中。

\textbf{13.} $3.3 \times 0.030 = \qty{0.10}{g}$ 至多。

\textbf{14.} $\eta = 2.7 / 3.91 = 0.69$：\qty{69}{\percent}。

\textbf{15.} 熔点尖锐，且与纯阿司匹林的熔点相同：
晶体为阿司匹林，且为纯品。

\textbf{16.} 熔化温度低于 \qty{135}{\celsius}，且熔程较宽：
粗产物不纯。

\textbf{17.} 粗产物含有阿司匹林和未反应的水杨酸；
重结晶已除去水杨酸。

\textbf{18.} $R_f = 3.1 / 5.0 = 0.62$。

\textbf{19.} 两步损失大致相同：从理论上可得的 \qty{3.91}{g}
到 \qty{3.3}{g} 粗产物，再经重结晶由 \qty{3.3}{g} 到 \qty{2.7}{g}，
每步约损失 \qty{0.6}{g}。由于粗产物
并非纯阿司匹林，因此第一步实际损失的
阿司匹林比 \qty{0.6}{g} 还多一些。

\textbf{20.} 合成的产率约为 \qty{69}{\percent}。
\end{solution}
